% =====================================================
% CHAPITRE 1 : PRÉSENTATION DE L'ORGANISME D'ACCUEIL
% =====================================================
\chapter{Présentation de l'organisme d'accueil}

% ---------------------------------------------------
\section*{Introduction}
\addcontentsline{toc}{section}{Introduction}
% ---------------------------------------------------

Ce premier chapitre pose le cadre institutionnel et professionnel dans lequel s'inscrit
ce projet de fin d'études. Avant de décrire les choix techniques et les résultats obtenus,
il est indispensable de comprendre l'environnement au sein duquel la problématique a émergé.
Nous présenterons successivement le réseau PwC à l'échelle internationale, son implantation
tunisienne, puis le département Risk Assurance Services et sa cellule Audit IT --- contexte
direct de production des données utilisées dans ce travail. Nous conclurons par la description
du cadre du stage et des missions confiées.

% ---------------------------------------------------
\section{PwC à l'échelle internationale}
% ---------------------------------------------------

\subsection{Historique et positionnement}

PwC est un réseau mondial de cabinets de services professionnels couvrant l'audit, le conseil,
la fiscalité et les transactions. Sa clientèle est très diversifiée : grands groupes
internationaux cotés, PME en pleine croissance, institutions publiques et organisations à but
non lucratif. Ce qui caractérise le réseau, c'est sa capacité à articuler une expertise
internationale uniformisée avec une connaissance fine des réalités locales --- une posture
particulièrement pertinente sur des marchés comme la Tunisie, où les spécificités
réglementaires et économiques sont fortes.

PwC est né en \textbf{1998} de la fusion de deux cabinets alors centenaires :
\textit{Price Waterhouse}, fondé à Londres en 1849, et \textit{Coopers \& Lybrand},
dont les origines remontent à 1854. L'entité fusionnée, d'abord connue sous le nom
PricewaterhouseCoopers, a adopté en septembre 2010 son identité actuelle abrégée « PwC ».
Aujourd'hui présent dans plus de \textbf{187 pays} avec \textbf{771 bureaux} et
\textbf{364~000 collaborateurs}, le réseau se positionne comme le premier cabinet mondial
d'audit et de conseil, devant Deloitte, EY et KPMG.

\subsection{PwC dans le Big Four}

PwC s'inscrit dans le cercle restreint des quatre grands cabinets mondiaux d'audit et de
conseil --- le \og Big Four \fg{} ---, aux côtés de Deloitte, Ernst \& Young et KPMG.
Chacun de ces réseaux repose sur le même modèle organisationnel : des entités nationales
juridiquement indépendantes, fédérées sous une marque commune et soumises à des standards
partagés en matière de qualité et d'éthique professionnelle.

\begin{table}[H]
\centering
\caption{Fiche signalétique de PwC International}
\label{tab:fiche_pwc_intl}
\begin{tabular}{ll}
\toprule
\textbf{Critère}              & \textbf{Information} \\
\midrule
Dénomination                  & PricewaterhouseCoopers (PwC) \\
Date de création              & 1998 (fusion Price Waterhouse \& Coopers \& Lybrand) \\
Siège social                  & Londres, Royaume-Uni \\
Secteur d'activité            & Audit, Conseil, Fiscalité, Transactions \\
Nombre de pays                & Plus de 187 \\
Nombre de bureaux             & Plus de 771 \\
Effectif mondial              & Plus de 364~000 collaborateurs \\
Appartenance                  & Big Four de l'audit \\
\bottomrule
\end{tabular}
\end{table}

\subsection{Statut juridique}

Sur le plan juridique, PwC s'appuie sur \textit{PricewaterhouseCoopers International
Limited} (PwCIL), une société de droit anglais qui joue un rôle de coordination
stratégique sans être elle-même prestataire de services. Chaque cabinet national est
une entité indépendante, ce qui lui permet de satisfaire aux exigences des pays imposant
un actionnariat local pour les cabinets d'audit. En contrepartie de l'usage de la marque
PwC, ces entités s'engagent à respecter les standards de qualité, de gestion du risque
et d'éthique professionnelle définis par PwCIL.

\subsection{Stratégie mondiale : « The New Equation »}

Pour anticiper les évolutions du marché --- notamment l'essor de l'intelligence
artificielle et les nouvelles exigences réglementaires ---, PwC a structuré sa stratégie
globale autour du programme \textbf{« The New Equation »}, articulé autour de deux axes :

\begin{itemize}
  \item \textbf{Building Trust} : répondre aux attentes en matière de transparence en
    combinant expertise d'audit, de fiscalité et de cybersécurité, enrichies par
    l'intelligence artificielle.
  \item \textbf{Sustained Outcomes} : mobiliser l'ensemble des expertises du réseau pour
    accompagner ses clients vers des résultats durables, couvrant stratégie, services
    numériques, fiscalité et conformité réglementaire.
\end{itemize}

% ---------------------------------------------------
\section{PwC Tunisie}
% ---------------------------------------------------

\subsection{Présentation générale}

Fondée en \textbf{1983}, PwC Tunisie s'impose comme l'une des premières sociétés
tunisiennes dans le domaine de l'audit et du conseil, forte de plus de quarante années
d'expérience au service du tissu économique local. Elle propose une gamme étendue de
services adaptés au marché tunisien, couvrant aussi bien les besoins des institutions
publiques que ceux des entreprises privées, des ONG et des bailleurs de fonds
internationaux.

\begin{table}[H]
\centering
\caption{Fiche signalétique de PwC Tunisie}
\label{tab:fiche_pwc_tunisie}
\begin{tabular}{ll}
\toprule
\textbf{Critère}              & \textbf{Information} \\
\midrule
Dénomination                  & PricewaterhouseCoopers Tunisie \\
Date de création              & 1983 \\
Siège social                  & Tunis, Tunisie \\
Secteur d'activité            & Audit, Conseil, Fiscalité, Transactions \\
Domaines d'intervention       & Secteur public, Privé, ONG, Bailleurs de fonds \\
Appartenance                  & Réseau PwC International \\
\bottomrule
\end{tabular}
\end{table}

\subsection{Offres de services de PwC Tunisie}

PwC Tunisie s'organise en quatre lignes de services complémentaires :

\paragraph{Audit.}
Évaluation du contrôle interne, conception d'organigrammes, développement de manuels de
procédures, construction de matrices de risques et de contrôles (Risk \& Control
Matrix --- RCM).

\paragraph{Transactions (Deals).}
Due diligence acheteur et vendeur, études de faisabilité, élaboration de business plans,
évaluation d'entreprises et gestion des processus d'ouverture de capital.

\paragraph{Services Juridiques et Fiscaux.}
Conseil fiscal, tax due diligence, audit fiscal, tax compliance et assistance lors des
contrôles fiscaux.

\paragraph{Consulting.}
Définition de stratégies de développement, refonte commerciale, amélioration de
l'efficacité opérationnelle et conseil en marketing.

\subsection{Références clients}

PwC Tunisie opère dans tous les domaines d'activité, avec une expertise forte dans les
services financiers, le secteur public, l'énergie, les projets d'infrastructure et les
télécommunications.

% ---------------------------------------------------
\section{Le département Risk Assurance Services}
% ---------------------------------------------------

\subsection{Présentation générale}

Le département \textbf{Risk Assurance Services} (anciennement Conseil Audit Formation) est
le pôle au sein duquel j'ai réalisé ce stage. Sa mission centrale : aider les organisations
clientes à identifier et cartographier leurs risques opérationnels et systémiques, puis à
mettre en place les dispositifs de contrôle appropriés. Concrètement, les missions couvrent
un spectre large --- de l'évaluation du contrôle interne à la revue des systèmes
d'information, en passant par l'assistance à la conception de procédures et la formation des
équipes audit.

\subsection{Offres de prestations}

Le département apporte une valeur ajoutée à travers trois formes d'intervention
complémentaires :

\paragraph{Outsourcing des travaux.}
Cette prestation s'adresse aux organisations ne disposant pas d'une fonction d'audit interne
constituée. PwC prend alors en charge l'intégralité des travaux d'audit, assurant une
couverture complète de l'ensemble des missions inscrites au plan d'audit interne.

\paragraph{Co-sourcing des travaux.}
L'équipe d'audit interne du client fait appel à l'expertise de PwC pour renforcer ses
capacités. Cette modalité collaborative favorise le transfert de connaissances et permet à
l'équipe interne de monter progressivement en compétences.

\paragraph{Services de conseil.}
PwC accompagne les entreprises dans l'évaluation de leurs opérations et dans la structuration
ou l'amélioration de leur fonction d'audit interne, afin de la rendre plus performante et
mieux alignée sur leurs enjeux stratégiques.

\subsection{La cellule Audit IT}

Au cœur du département Risk Assurance Services, la cellule \textbf{Audit IT} se consacre
à l'examen des systèmes d'information qui alimentent la comptabilité et les états financiers
des clients : ERP, systèmes de gestion de trésorerie, plateformes de paiement. Son rôle
est de vérifier que ces environnements produisent des données intègres, fiables et conformes
aux normes d'audit --- ce qui conditionne directement la validité des assertions financières
sur lesquelles repose l'opinion du commissaire aux comptes.

C'est précisément dans ce cadre opérationnel que les données de ce projet ont été collectées :
un corpus de transactions financières extraites du système d'information d'un client de PwC
Tunisie, dans le cadre d'une mission d'audit IT active. Pour des raisons de confidentialité,
toutes les informations permettant d'identifier les parties ont été anonymisées.

% ---------------------------------------------------
\section{Cadre du stage et projet de fin d'études}
% ---------------------------------------------------

\subsection{Objectifs du stage}

Ce stage de fin d'études, réalisé au sein du département Risk Assurance Services de PwC
Tunisie, a pour objectif principal de concevoir et de déployer un \textbf{système de
détection intelligente d'anomalies financières} fondé sur des techniques avancées
d'apprentissage automatique. Il s'inscrit dans une démarche d'innovation visant à renforcer
les capacités analytiques de l'équipe d'audit, en automatisant partiellement l'identification
des transactions suspectes et en produisant des explications compréhensibles par les
analystes métier.

\subsection{Missions confiées}

Les missions confiées dans le cadre de ce stage couvrent l'intégralité du cycle CRISP-DM :

\begin{itemize}
  \item Compréhension du contexte métier de la détection de fraude financière en
    environnement d'audit ;
  \item Analyse exploratoire du corpus de transactions réelles
    ($\sim$200~000 opérations, ratio fraude 1:774) ;
  \item Conception et implémentation d'un pipeline de préparation des données
    (nettoyage, ingénierie de variables, gestion du déséquilibre via SMOTE) ;
  \item Développement et comparaison de trois approches de modélisation : régression
    logistique, forêt aléatoire et autoencoder non supervisé ;
  \item Mise en place d'une couche d'explicabilité double : méthodes SHAP et LIME, et
    génération d'explications en langage naturel via \textbf{Ollama} ;
  \item Déploiement d'un pipeline reproductible via le script d'orchestration
    \texttt{run\_all.py}, avec couverture de tests unitaires.
\end{itemize}

\subsection{Présentation du projet}

Le projet vise à construire une \textbf{chaîne de traitement complète et reproductible} pour
la détection de transactions frauduleuses. Conduit selon la méthodologie \textbf{CRISP-DM},
il répond à la problématique suivante : \textit{comment détecter efficacement les fraudes
dans un corpus fortement déséquilibré, tout en produisant des explications compréhensibles
pour les auditeurs ?}

La solution développée repose sur trois contributions complémentaires :

\begin{itemize}
  \item \textbf{Modélisation multi-approches} : comparaison d'un modèle linéaire, d'un
    modèle ensembliste supervisé et d'un autoencoder profond non supervisé, évalués sur
    des métriques adaptées au déséquilibre (F1-score, PR-AUC, rappel) ;
  \item \textbf{Gestion du déséquilibre} : stratégies de pondération des classes et de
    rééchantillonnage SMOTE pour traiter le ratio fraude/légitime de 1:774 ;
  \item \textbf{Explicabilité des décisions} : identification des variables les plus
    contributives via SHAP et LIME, et génération d'explications en langage naturel à
    destination des analystes métier via le LLM Ollama.
\end{itemize}

% ---------------------------------------------------
\section*{Conclusion}
\addcontentsline{toc}{section}{Conclusion}
% ---------------------------------------------------

Ce premier chapitre a établi le cadre institutionnel et organisationnel dans lequel s'inscrit
ce projet de fin d'études. À travers la présentation de PwC en tant que réseau international
de référence, puis de PwC Tunisie et de son département Risk Assurance Services, nous avons
mis en lumière l'environnement professionnel au sein duquel ce stage a été effectué, ainsi
que le rôle stratégique de la cellule Audit IT. La problématique identifiée --- détecter
automatiquement les transactions frauduleuses dans un corpus réel fortement déséquilibré,
tout en garantissant l'interprétabilité des résultats --- et les missions confiées ont posé
les bases nécessaires à la compréhension des enjeux du projet.

Le chapitre suivant sera consacré à la compréhension métier, première phase de la
méthodologie CRISP-DM, dans laquelle nous formaliserons la problématique, les objectifs
et les critères de succès du projet.
